\problem{18}{\thtxt{(สอวน ม.4 2551) ค่าคงตัวหลักมูลสามค่าในฟิสิกส์คือ ค่าคงตัวโน้มถ่วงสากล} $G=6.7\times10^{-11}\,\mathrm{m^3\,kg^{-1}\,s^{-2}}$ \thtxt{อัตราเร็วของแสงในสูญญากาศ} $c=3.0\times10^8\,\mathrm{m\,s^{-1}}$ \thtxt{และค่าคงตัวของพลังค์} $h=6.6\times10^{-34}\,\mathrm{J\,s}$ \thtxt{ถ้าเราผสมค่าคงตัวสามค่านี้เหมาะ ๆ เช่น เอากำลังต่าง ๆ ของปริมาณเหล่านี้คูณหารกัน เราจะได้ปริมาณที่มีหน่วยของเวลา เราเรียกเวลาที่ได้นี้ว่า เวลาพลังค์ นักฟิสิกส์คิดว่าเวลานี้เป็นอายุของเอกภพก่อนที่กฎฟิสิกส์ต่าง ๆ ที่เรารู้จักในปัจจุบันจะใช้ได้ ให้ตอบสูตรของเวลาพลังค์นี้ในรูปของค่าคงตัวหลักมูลสามตัวที่กำหนดให้ และให้หาว่า เวลานี้มีขนาดเท่าใดในหน่วยวินาที (แนะให้เขียนหน่วยปริมาณต่าง ๆ ในรูปของหน่วยหลักมูล เมตร กิโลกรัม และวินาที ถ้าปริมาณสองปริมาณเท่ากัน หน่วยของปริมาณทั้งสองนั้นในรูปของหน่วยหลักมูล เมตร กิโลกรัม และวินาที จะต้องเท่ากัน)}}
